\section{Generic Algebra for a Second-Order Accounting Wedge}

This appendix records the generic algebra behind the wedge construction. The
important sequencing point is that the Taylor expansion is applied to the
nonlinear equilibrium condition, not to an already log-linearized equation. Let
an equilibrium condition be written as
\begin{equation}
    \E_t
    \left[
        F(z_t,z_{t+1})
    \right]
    =
    0,
    \label{eq:generic_equilibrium_condition}
\end{equation}
where \(z_t\) stacks local coordinates for the relevant endogenous and exogenous
variables. These coordinates can be log deviations, such as
\(Z_t=\bar{Z}\exp(z_t)\), but this is only a change of variables. It is not a
first-order approximation. Let \(\bar{z}=0\) denote the deterministic steady
state in those coordinates and define \(\tilde{z}_t=z_t-\bar{z}\). A
second-order expansion gives
\begin{equation}
    0
    \approx
    F(\bar{z},\bar{z})
    +
    F_0 \tilde{z}_t
    +
    F_1 \E_t \tilde{z}_{t+1}
    +
    \frac{1}{2}
    \E_t
    \left[
        \tilde{z}_{t+1}'
        H_1
        \tilde{z}_{t+1}
    \right]
    +
    \text{terms involving } \tilde{z}_t^2.
    \label{eq:generic_second_order_expansion}
\end{equation}
At the deterministic steady state, \(F(\bar{z},\bar{z})=0\). The conditional
quadratic term can be decomposed as
\begin{equation}
    \E_t
    \left[
        \tilde{z}_{t+1}'
        H_1
        \tilde{z}_{t+1}
    \right]
    =
    (\E_t\tilde{z}_{t+1})'
    H_1
    (\E_t\tilde{z}_{t+1})
    +
    \tr(H_1\Omega_t),
    \label{eq:mean_variance_decomposition}
\end{equation}
where
\begin{equation}
    \Omega_t
    =
    \Var_t(\tilde{z}_{t+1}).
    \label{eq:omega_covariance}
\end{equation}

For an accounting exercise focused on uncertainty, collect all terms involving
\(\Omega_t\) into a wedge:
\begin{equation}
    \omega^F_t
    \equiv
    -A_F^{-1}
    \left[
        \frac{1}{2}\tr(H_1\Omega_t)
    \right],
    \label{eq:generic_wedge}
\end{equation}
where \(A_F\) is the coefficient on the endogenous variable being isolated.

If one first replaces Equation \eqref{eq:generic_equilibrium_condition} by its
first-order approximation,
\begin{equation}
    F_0 \tilde{z}_t
    +
    F_1 \E_t\tilde{z}_{t+1}
    =
    0,
    \label{eq:first_order_discard}
\end{equation}
then the Hessian term \(H_1\) has already been discarded. In that case there is
no \(\tr(H_1\Omega_t)\) term left to interpret as an uncertainty wedge.

For the capital Euler condition, the generic wedge in Equation
\eqref{eq:generic_wedge} corresponds to the log investment wedge in Equation
\eqref{eq:investment_uncertainty_wedge}. This is why the investment wedge is the
canonical BCA wedge most closely connected to second-moment risk corrections.
